% Wersja źródłowa: audyt II, jednostka 27
\setcounter{section}{26}
\noindent\textbf{Wyniki wykorzystywane.}
Przyjmujemy znane twierdzenia Hahna--Banacha,
Banacha--Steinhausa, o odwzorowaniu otwartym,
twierdzenie Eberleina--\v Smuliana oraz podstawy
ortogonalności i bazy ortonormalnej przestrzeni Hilberta.
\par\medskip
\section{Bazy Schaudera, układy biortogonalne i zbieżność bezwarunkowa}
\label{sec:bazy27}

W tym rozdziale $X$ oznacza przestrzeń Banacha nad $\K\in\{\R,\C\}$,
$X'$ jej przestrzeń dualną topologiczną, a $\mathcal L(X,Y)$
przestrzeń operatorów liniowych ciągłych. Bazy ortonormalne przestrzeni
Hilberta omówiono we wcześniejszym rozdziale o przestrzeniach Hilberta; tutaj badamy bazy
\emph{uporządkowane}, właściwe ogólnym przestrzeniom Banacha.
Nie należy mylić baz Schaudera z bazami Hamela: w przypadku tych
ostatnich dopuszcza się wyłącznie \emph{skończone} kombinacje liniowe.

\subsection{Definicja, współczynniki i ciągłość rzutów częściowych}

\begin{definicja}[Baza Schaudera i współczynniki]
\label{def:schauder27}
Niech $X$ będzie przestrzenią Banacha, a $(e_n)_{n\ge1}\subset X$
ciągiem wektorów niezerowych. Nazywamy go \emph{bazą Schaudera}
przestrzeni $X$, jeśli dla każdego $x\in X$ istnieje dokładnie jeden
ciąg skalarów $(a_n)_{n\ge1}$, dla którego
$x=\lim_{N\to\infty}\sum_{n=1}^N a_ne_n$ w normie $X$.
Skalar $a_n$ oznaczamy $e_n^*(x)$ i nazywamy $n$-tym
\emph{współczynnikiem} wektora $x$. Odwzorowanie
$P_Nx=\sum_{n=1}^Ne_n^*(x)e_n$ nazywamy
\emph{rzutem częściowym}; kładziemy też $P_0=0$.
\end{definicja}

Jednoznaczność implikuje liniowość $e_n^*$ i $P_N$; nie zakładaliśmy
natomiast jeszcze ich ciągłości. Jest to istotne: ciągłość
współczynników stanowi treść następnego twierdzenia.

\begin{lemat}[Przestrzeń ciągów współczynników]
\label{lem:coefficient-space27}
Niech $X$ będzie przestrzenią Banacha, a $(e_n)_{n\ge1}\subset X$
ciągiem wektorów niezerowych. Zbiór
\[
E_e=\bigl\{a=(a_n)\in\K^{\N}:\sum_{n=1}^{\infty}a_ne_n
\text{ jest zbieżny w }X\bigr\}
\]
z normą $\|a\|_{E_e}=\sup_{N\ge1}\|\sum_{n=1}^Na_ne_n\|_X$
jest przestrzenią Banacha. Ponadto odwzorowanie
$S:E_e\to X$, $S(a)=\sum_{n\ge1}a_ne_n$, jest liniowe
ciągłe i spełnia $\|S\|\le1$.
\end{lemat}
\begin{proof}
Norma jest skończona dzięki zbieżności sum częściowych.
Jeżeli wszystkie sumy częściowe znikają, to z różnic kolejnych
sum i z $e_n\ne0$ wynika $a_n=0$ dla wszystkich $n$.
Niech $(a^{(k)})$ będzie ciągiem Cauchy'ego w $E_e$.
Dla każdego $n$ mamy
$|a_n^{(k)}-a_n^{(\ell)}|\,\|e_n\|
\le2\|a^{(k)}-a^{(\ell)}\|_{E_e}$,
więc istnieje granica $a_n=\lim_k a_n^{(k)}$.
Przechodząc do granicy przy ustalonym $N$, otrzymujemy
$\|\sum_{n=1}^N(a_n^{(k)}-a_n)e_n\|\le\eta$
jednostajnie w $N$, gdy $k$ jest dostatecznie duże,
bo analogiczne oszacowanie z $a^{(\ell)}$ jest
prawdziwe dla każdego $\ell$.
Wybieramy takie $k$. Dla $M>N$ mamy
\[
\Big\|\sum_{n=N+1}^Ma_ne_n\Big\|
\le2\eta+\Big\|\sum_{n=N+1}^Ma_n^{(k)}e_n\Big\|.
\]
W ostatnim wyrażeniu, dla ustalonego $k$, sumy ogonowe
są dowolnie małe przy dużym $N$. Ponieważ $\eta>0$
można przyjąć dowolnie małe, szereg dla $a$ jest zbieżny.
Ponadto $\|a^{(k)}-a\|_{E_e}\to0$.
Dla każdego $a\in E_e$ mamy
$\|S(a)\|\le\|a\|_{E_e}$, co dowodzi tezy o $S$.
\end{proof}

\begin{twierdzenie}[Ciągłość współczynników i stała bazy]
\label{thm:basis-continuity27}
Jeżeli $(e_n)_{n\ge1}$ jest bazą Schaudera przestrzeni
Banacha $X$, to $e_n^*\in X'$ dla każdego $n$, a każdy
$P_N\in\mathcal L(X,X)$. Ponadto
\[
K_e:=\sup_{N\ge1}\|P_N\|<\infty,\qquad K_e\ge1.
\]
Liczbę $K_e$ nazywamy \emph{stałą bazy}.
Dla każdego $n\ge1$ zachodzi $e_n^*(e_m)=\delta_{nm}$
oraz $\|e_n^*\|\,\|e_n\|\le2K_e$.
\end{twierdzenie}
\begin{proof}
Dla bazy Schaudera odwzorowanie $S:E_e\to X$
z Lematu~\ref{lem:coefficient-space27} jest liniową ciągłą
bijekcją. Twierdzenie o odwzorowaniu otwartym zapewnia
ciągłość $S^{-1}$. Dla każdego $N$ zachodzi
$\|P_Nx\|=\|\sum_{n=1}^Ne_n^*(x)e_n\|
\le\|S^{-1}x\|_{E_e}$,
stąd $K_e\le\|S^{-1}\|<\infty$.
Dla $n\ge1$ tożsamość
$(P_n-P_{n-1})x=e_n^*(x)e_n$
daje $\|e_n^*\|\,\|e_n\|\le\|P_n\|+\|P_{n-1}\|
\le2K_e$. Zatem wszystkie współczynniki są ciągłe.
Biortogonalność wynika z jednoznaczności reprezentacji $e_m$.
Każdy $P_N$ jest niezerowym rzutem (pozostawia $e_1$),
a więc $\|P_N\|\ge1$.
\end{proof}

\begin{wniosek}[Normalizacja i ośrodkowość]
\label{cor:separable27}
Każda przestrzeń Banacha z bazą Schaudera jest ośrodkowa.
Jeżeli $(e_n)$ jest bazą Schaudera, to
$(e_n/\|e_n\|)$ również jest bazą Schaudera o tej samej stałej
bazy; jej współczynniki są równe $\|e_n\|e_n^*(x)$.
\end{wniosek}
\begin{proof}
Skończone kombinacje $e_n$ o współczynnikach w $\mathbb Q$
(dla $\K=\C$ w $\mathbb Q+i\mathbb Q$) tworzą zbiór przeliczalny
i gęsty w $X$, bo sumy częściowe bazy przybliżają każdy $x$.
Zmiana normalizacji nie zmienia ani sum częściowych,
ani ich granic, zatem nie zmienia operatorów $P_N$.
\end{proof}

\begin{definicja}[Układ biortogonalny i ciąg bazowy]
\label{def:biorthogonal27}
Niech $X$ będzie przestrzenią Banacha.
Układ $(e_n,f_n)_{n\ge1}\subset X\times X'$ nazywamy
\emph{biortogonalnym}, jeżeli $f_n(e_m)=\delta_{nm}$
dla wszystkich $m,n$. Ciąg $(e_n)$ nazywamy
\emph{ciągiem bazowym}, jeżeli jest bazą Schaudera swojej
normowo domkniętej otoczki liniowej $[e_n]
=\overline{\operatorname{span}\{e_n:n\ge1\}}^{\,\|\cdot\|}$.
\end{definicja}

\begin{twierdzenie}[Kryterium bazy przez rzuty]
\label{thm:biorthogonal-criterion27}
Niech $X$ będzie przestrzenią Banacha, zaś
$(e_n,f_n)_{n\ge1}\subset X\times X'$ układem biortogonalnym,
dla którego $\operatorname{span}\{e_n:n\ge1\}$ jest gęsta w $X$.
Oznaczmy $P_Nx=\sum_{n=1}^Nf_n(x)e_n$.
Wówczas $(e_n)$ jest bazą Schaudera $X$, której
funkcjonałami współrzędnych są $f_n$, wtedy i tylko wtedy,
gdy $\sup_N\|P_N\|<\infty$.
\end{twierdzenie}
\begin{proof}
Warunek jest konieczny na mocy
Twierdzenia~\ref{thm:basis-continuity27}.
W przeciwną stronę niech $K=\sup_N\|P_N\|<\infty$.
Jeśli $y$ jest skończoną kombinacją $e_n$, to $P_Ny=y$
dla dostatecznie dużych $N$. Dla dowolnego $x\in X$
wybieramy takie $y$ bliskie $x$; wówczas
$\|x-P_Nx\|\le(1+K)\|x-y\|$ dla dużych $N$.
Zatem $P_Nx\to x$ w normie.
Jeżeli również $x=\sum_na_ne_n$, to przez ciągłość $f_j$
i biortogonalność dostajemy $f_j(x)=a_j$.
To dowodzi jednoznaczności rozwinięcia.
\end{proof}

\begin{wniosek}[Kryterium ciągu bazowego]
\label{cor:basic-sequence27}
Niech $(e_n)_{n\ge1}$ będzie ciągiem wektorów niezerowych
w przestrzeni Banacha $X$. Jest to ciąg bazowy wtedy i tylko wtedy,
gdy jest liniowo niezależny i istnieje $K<\infty$ takie, że
$\|\sum_{n=1}^Na_ne_n\|
\le K\|\sum_{n=1}^Ma_ne_n\|$
dla wszystkich $M\ge N\ge1$ i wszystkich $a_1,\ldots,a_M\in\K$.
\end{wniosek}
\begin{proof}
Dla ciągu bazowego nierówność wynika z ograniczoności
rzutów częściowych na $[e_n]$.
Odwrotnie, na algebraicznej otoczce liniowej $e_n$
jednoznacznie definiujemy $P_N$ przez ucięcie współczynników.
Założona nierówność daje $\|P_N\|\le K$.
Każdy $P_N$ przedłuża się do $[e_n]$ przez ciągłość;
wówczas $P_Nx\to x$ przez gęstość skończonych kombinacji.
Współczynniki są ciągłe, ponieważ
$(P_n-P_{n-1})x$ należy do jednowymiarowej przestrzeni
$\K e_n$, a ciągłość funkcjonałów współrzędnych
wynika z podzielenia przez stały niezerowy wektor $e_n$.
Stosujemy Twierdzenie~\ref{thm:biorthogonal-criterion27}
do przestrzeni $[e_n]$.
\end{proof}

\subsection{Bazy monotoniczne i klasyczne przykłady}

\begin{definicja}[Baza monotoniczna]
\label{def:monotonic27}
Bazę Schaudera $(e_n)$ nazywamy \emph{monotoniczną}, gdy
$\|P_N\|=1$ dla wszystkich $N\ge1$.
\end{definicja}

\begin{przyklad}[Przestrzenie ciągowe]
\label{ex:sequence-bases27}
W przestrzeniach Banacha $c_0(\N,\K)$ i $\ell^p(\N,\K)$
dla $1\le p<\infty$ wektory standardowe
$e_n=(\delta_{nj})_{j\ge1}$ tworzą bazę Schaudera.
Współczynnik $e_n^*(x)$ to $n$-ta współrzędna,
a $P_N(x)$ zeruje wszystkie współrzędne po $N$.
Dla $\ell^p$ zachodzi
$\|x-P_Nx\|_p^p=\sum_{j>N}|x_j|^p\to0$,
a dla $c_0$ mamy $\|x-P_Nx\|_\infty=\sup_{j>N}|x_j|\to0$.
Stąd $\|P_N\|=1$ i bazy są monotoniczne.
W $\ell^\infty$ ciąg standardowy nie jest bazą Schaudera:
jego domknięta otoczka liniowa równa jest $c_0\ne\ell^\infty$.
\end{przyklad}

\begin{przyklad}[Baza sumująca w $c_0$]
\label{ex:summing-basis27}
Niech $s_n=\sum_{j=1}^ne_j\in c_0$.
Dla $x=(x_j)\in c_0$ przyjmijmy $a_n=x_n-x_{n+1}$.
Wtedy $\sum_{n=1}^Na_ns_n$ ma na pozycjach $j\le N$
wartość $x_j-x_{N+1}$, a po pozycjach $j>N$ wartość $0$.
Zatem
$\|x-\sum_{n=1}^Na_ns_n\|_\infty
\le2\sup_{j>N}|x_j|\to0$.
Jednoznaczność współczynników wynika z tożsamości
$a_n=x_n-x_{n+1}$; układ $(s_n)$ jest więc bazą Schaudera.
Jej rzuty częściowe mają normę $2$ dla każdego $N\ge1$,
ponieważ $\|P_Nx\|_\infty
\le\|x\|_\infty+|x_{N+1}|\le2\|x\|_\infty$,
a równość otrzymujemy przez wybór $x_1=1$, $x_{N+1}=-1$
i pozostałych współrzędnych równych zeru (dla $N\ge1$).
Ta baza nie jest monotoniczna.
\end{przyklad}

\subsection{Bazy bezwarunkowe i zmiana kolejności sumowania}

\begin{definicja}[Baza bezwarunkowa]
\label{def:unconditional27}
Bazę Schaudera $(e_n)$ przestrzeni Banacha $X$ nazywamy
\emph{bezwarunkową}, jeżeli dla każdego $x\in X$ i każdej
permutacji $\pi:\N\to\N$ szereg
$\sum_{n=1}^\infty e_{\pi(n)}^*(x)e_{\pi(n)}$
jest zbieżny w normie $X$ do $x$.
Dla skończonego $F\subseteq\N$ piszemy
$P_Fx=\sum_{n\in F}e_n^*(x)e_n$.
\end{definicja}

\begin{lemat}[Permutacje a sumy po zbiorach skończonych]
\label{lem:net-unconditional27}
Niech $(u_n)$ będzie ciągiem w przestrzeni Banacha $X$
i załóżmy, że szereg $\sum_nu_n$ jest zbieżny w zadanej
kolejności do $u\in X$. Następujące warunki są równoważne:
\begin{enumerate}[label=\textup{(\roman*)}]
\item wszystkie permutacje szeregu są zbieżne do $u$;
\item sieć $(\sum_{n\in F}u_n)_{F\Subset\N}$, skierowana przez inkluzję,
 zbiega do $u$;
\item dla każdego $\varepsilon>0$ istnieje skończony $F_0\subseteq\N$
 taki, że $\|\sum_{n\in F}u_n\|<\varepsilon$
 dla każdego skończonego $F\subseteq\N\setminus F_0$.
\end{enumerate}
\end{lemat}
\begin{proof}
Jeśli (iii) zachodzi, to dla skończonych
$F,G\supseteq F_0$ różnica ich sum jest różnicą sum po
$F\setminus G$ i $G\setminus F$. Oba zbiory są rozłączne
z $F_0$, zatem nierówność trójkąta daje normę różnicy
mniejszą niż $2\varepsilon$. Zatem sieć jest
siecią Cauchy'ego. W przestrzeni Banacha ma granicę,
która musi być równa $u$, ponieważ zawiera ona jako podsieć
sumy po $\{1,\ldots,N\}$.
Implikacja (ii) $\Rightarrow$ (i) wynika stąd, że
zbiory $\{\pi(1),\ldots,\pi(N)\}$ wyczerpują $\N$.
Załóżmy, że (iii) nie zachodzi. Istnieje $\varepsilon>0$
takie, że po usunięciu dowolnego zbioru skończonego
można znaleźć zbiór skończony $F$ z normą sumy
przynajmniej $\varepsilon$. Rekurencyjnie konstruujemy
rozłączne zbiory skończone $F_k$, których elementy są coraz
większe, i których sumy mają normę przynajmniej $\varepsilon$.
Każde $F_k$ wybieramy poza poprzednimi zbiorami i poza
$\{1,\ldots,k\}$. Permutację $\N$ budujemy blokami,
wyliczając kolejno elementy $F_k$ i wstawiając
między blokami najmniejszy dotąd niewykorzystany indeks.
Wyrazy tak przestawionego szeregu mają sumy częściowe,
które na końcach każdego bloku $F_k$ różnią się o co najmniej
$\varepsilon$. Szereg po tej permutacji nie jest więc
zbieżny, co przeczy (i).
\end{proof}

\begin{twierdzenie}[Kryterium bezwarunkowości przez rzuty]
\label{thm:unconditional-projections27}
Niech $(e_n)$ będzie bazą Schaudera przestrzeni Banacha $X$.
Wówczas następujące warunki są równoważne:
\begin{enumerate}[label=\textup{(\roman*)}]
\item $(e_n)$ jest bazą bezwarunkową;
\item istnieje $U<\infty$ takie, że $\|P_F\|\le U$
 dla każdego skończonego $F\subseteq\N$;
\item dla każdego $x\in X$ rodzina
$\{P_Fx:F\Subset\N\}$ jest ograniczona w normie.
\end{enumerate}
Najmniejsze $U$ z (ii) nazywamy \emph{stałą bezwarunkowości
przez opuszczanie współczynników}.
\end{twierdzenie}
\begin{proof}
Z (i) i Lematu~\ref{lem:net-unconditional27} wynika, że
$P_Fx\to x$ w sieci zbiorów skończonych, zatem rodzina
$\{P_Fx\}$ jest ograniczona. To daje (iii).
Twierdzenie Banacha--Steinhausa zastosowane do operatorów
$P_F\in\mathcal L(X,X)$ daje (iii) $\Rightarrow$ (ii).
Jeżeli zachodzi (ii), to dla $x\in X$ i dla przybliżenia
$y\in\operatorname{span}\{e_n\}$ mamy
$\|P_F(x-y)\|\le U\|x-y\|$.
Gdy $F$ zawiera wszystkie indeksy występujące w $y$,
to $P_Fy=y$, skąd $\|P_Fx-x\|\le(U+1)\|x-y\|$.
Gęstość skończonych kombinacji implikuje zbieżność
sieci $P_Fx\to x$, a Lemat~\ref{lem:net-unconditional27}
daje (i).
\end{proof}

\begin{wniosek}[Bazy ortonormalne i bezwarunkowość]
\label{cor:hilbert-unconditional27}
Niech $H$ będzie ośrodkową, nieskończenie wymiarową
przestrzenią Hilberta, a $(e_n)_{n\ge1}$ dowolną numeracją
jej bazy ortonormalnej. Wtedy $(e_n)$ jest monotoniczną
bazą Schaudera i bezwarunkową o stałej $U=1$.
\end{wniosek}
\begin{proof}
Twierdzenie Parsevala daje zbieżność rozwinięcia
$\sum_n\langle x,e_n\rangle e_n$ do $x$.
Dla każdego skończonego $F$ twierdzenie Pitagorasa
daje $\|P_Fx\|^2=\sum_{n\in F}|\langle x,e_n\rangle|^2
\le\|x\|^2$. Zarówno $P_F$, jak i $P_N$ są
niezerowymi rzutami, więc mają normę dokładnie $1$.
Twierdzenie~\ref{thm:unconditional-projections27} kończy dowód.
\end{proof}

\begin{przyklad}[Baza warunkowa w przestrzeni $c_0$]
\label{ex:conditional-summing27}
Baza sumująca $(s_n)$ z Przykładu~\ref{ex:summing-basis27}
nie jest bezwarunkowa. Weźmy $x\in c_0$ o współrzędnych
$x_n=(-1)^n/\sqrt n$. Wtedy
$a_n=s_n^*(x)=x_n-x_{n+1}=(-1)^n(1/\sqrt n+1/\sqrt{n+1})$.
Szereg skalarów $\sum_n a_n$ jest zbieżny (jego suma to
$x_1$), ale nie jest bezwzględnie zbieżny.
Gdyby szereg wektorowy $\sum_n a_ns_n$ był bezwarunkowo
zbieżny w $c_0$, to po zastosowaniu ciągłego
funkcjonału $\delta_1(y)=y_1$ uzyskalibyśmy bezwarunkową
zbieżność $\sum_n a_n$ w $\K$, co jest równoważne
zbieżności bezwzględnej. Otrzymujemy sprzeczność.
\end{przyklad}

\begin{przyklad}[Bezwarunkowość nie oznacza absolutnej zbieżności]
\label{ex:not-absolute27}
W przestrzeni Hilberta $\ell^2$ szereg
$\sum_{n\ge1}n^{-1}e_n$ jest zbieżny bezwarunkowo
(każdy wybór skończonej rodziny wyrazów ma kwadrat
normy równy sumie $n^{-2}$), ale
$\sum_{n\ge1}\|n^{-1}e_n\|=\sum_{n\ge1}n^{-1}=\infty$.
\end{przyklad}

\subsection{Bazy ściągające i zupełne w sensie ograniczonym}

\begin{definicja}[Dwa warunki na bazę]
\label{def:shrinking-bcomplete27}
Niech $(e_n)$ będzie bazą Schaudera przestrzeni Banacha $X$,
a $(e_n^*)\subset X'$ jej funkcjonałami współrzędnych.
Bazę nazywamy \emph{ściągającą} (ang. \emph{shrinking}), gdy
$\overline{\operatorname{span}\{e_n^*:n\ge1\}}^{\,\|\cdot\|}=X'$.
Nazywamy ją \emph{zupełną w sensie ograniczonym}
(ang. \emph{boundedly complete}), jeżeli dla każdego ciągu
skalarów $(a_n)$ warunek
$\sup_N\|\sum_{n=1}^Na_ne_n\|<\infty$
implikuje zbieżność szeregu $\sum_na_ne_n$ w normie $X$.
\end{definicja}

\begin{twierdzenie}[Kryterium refleksywności przez bazę]
\label{thm:james-basis27}
Niech $X$ będzie przestrzenią Banacha z bazą Schaudera $(e_n)$.
Wtedy $X$ jest refleksywna wtedy i tylko wtedy, gdy
$(e_n)$ jest jednocześnie ściągająca i zupełna w sensie
ograniczonym.
\end{twierdzenie}
\begin{proof}
Załóżmy najpierw refleksywność $X$.
Aby dowieść zupełności w sensie ograniczonym,
niech $s_N=\sum_{n=1}^Na_ne_n$ będzie ograniczony.
Kula zawierająca wszystkie $s_N$ jest słabo zwarta.
Istnieje zatem podciąg słabo zbieżny do pewnego $x\in X$
na mocy twierdzenia Eberleina--\v Smuliana.
Dla każdego $j$ ciąg $e_j^*(s_N)$ jest od pewnego miejsca
stały i równy $a_j$, więc $e_j^*(x)=a_j$.
Stąd $P_Nx=s_N$ dla wszystkich $N$, a definicja bazy
daje $s_N\to x$ w normie.

Niech teraz $f\in X'$. Oznaczmy
$Y_N=\overline{\operatorname{span}\{e_j:j>N\}}^{\,\|\cdot\|}$.
Pokażemy, że $\|f|_{Y_N}\|\to0$, gdzie norma oznacza
normę ograniczenia $f$ do domkniętej podprzestrzeni $Y_N$.
Gdyby tak nie było, istniałyby $\varepsilon>0$, liczby
$N_k\nearrow\infty$ i wektory $y_k\in Y_{N_k}$
o normie $1$ takie, że $|f(y_k)|\ge\varepsilon$.
Wektory te można wybierać rekurencyjnie, gdyż $Y_{N+1}\subseteq Y_N$,
a więc normy ograniczeń tworzą ciąg nierosnący.
Refleksywność i twierdzenie Eberleina--\v Smuliana
dają podciąg $(y_{k_\ell})$ słabo zbieżny do $y\in X$.
Dla każdego ustalonego $j$ mamy
$e_j^*(y_{k_\ell})=0$ dla dużych $\ell$, więc $e_j^*(y)=0$.
Jednoznaczność rozwinięcia wymusza $y=0$,
co jest sprzeczne z $|f(y_{k_\ell})|\ge\varepsilon$.
Ponieważ $(I-P_N)X\subseteq Y_N$ i $\|I-P_N\|\le1+K_e$,
dostajemy $\|f-P_N'f\|\le(1+K_e)\|f|_{Y_N}\|\to0$.
Funkcjonał $P_N'f=f\circ P_N$ jest skończoną kombinacją
$e_1^*,\ldots,e_N^*$; baza jest zatem ściągająca.

Odwrotnie, niech baza ma obie własności i niech
$\Phi\in X''$. Kładziemy $a_n=\Phi(e_n^*)$ oraz
$s_N=\sum_{n=1}^Na_ne_n$. Dla $f\in X'$ mamy
$f(s_N)=\Phi(P_N'f)$, a więc
$|f(s_N)|\le K_e\|\Phi\|\,\|f\|$.
Z definicji normy wynika $\|s_N\|\le K_e\|\Phi\|$.
Zupełność w sensie ograniczonym daje
$x=\lim_Ns_N\in X$. Dla każdego $n$ otrzymujemy
$e_n^*(x)=\Phi(e_n^*)$.
Zatem $J_Xx$ i $\Phi$ pokrywają się na liniowej otoczce
$\{e_n^*\}$, która jest gęsta w $X'$ dzięki ściąganiu bazy.
Ciągłość obu funkcjonałów na $X'$ daje $J_Xx=\Phi$.
Kanoniczne odwzorowanie $J_X:X\to X''$ jest surjektywne,
czyli $X$ jest refleksywna.
\end{proof}

\begin{przyklad}[Przykłady rozdzielające oba warunki]
\label{ex:james-examples27}
Standardowa baza przestrzeni $c_0$ jest ściągająca,
gdyż $c_0'=\ell^1$, a jej funkcjonały współrzędnych
odpowiadają wektorom standardowym w $\ell^1$.
Nie jest zupełna w sensie ograniczonym: sumy
$\sum_{n=1}^Ne_n$ mają normę $1$ w $c_0$, ale
nie tworzą ciągu Cauchy'ego.
Standardowa baza przestrzeni $\ell^1$ jest zupełna w sensie
ograniczonym, ponieważ ograniczoność sum częściowych oznacza
$\sup_N\sum_{n=1}^N|a_n|<\infty$.
Nie jest ściągająca, gdyż $(\ell^1)'=\ell^\infty$,
a domknięcie liniowej otoczki funkcjonałów współrzędnych
to właściwa podprzestrzeń $c_0$ przestrzeni $\ell^\infty$.
\end{przyklad}

\subsection{Czego istnienie bazy nie gwarantuje?}

\begin{uwaga}[Baza ortonormalna a baza Hamela]
\label{rem:hamel-vs-schauder27}
Jeżeli przestrzeń Hilberta jest ośrodkowa i nieskończenie
wymiarowa, to przeliczalna baza ortonormalna jest bazą
Schaudera, lecz nigdy bazą Hamela całej przestrzeni:
wektory z bazy Hamela dają skończone kombinacje liniowe.
Co więcej, dla nieskończenie wymiarowej przestrzeni Hilberta
każda baza Hamela jest nieprzeliczalna. Dowody tych faktów,
a także przeliczalność niezerowych współczynników
względem dowolnej, nawet nieprzeliczalnej rodziny
ortonormalnej, znajdują się we wcześniejszym rozdziale o przestrzeniach Hilberta.
\end{uwaga}

\begin{uwaga}[Ośrodkowość nie wystarcza]
\label{rem:enflo27}
Istnieją ośrodkowe przestrzenie Banacha bez bazy Schaudera.
Wynik ten ma charakter głęboki i nie wynika z samej
nierówności Banacha--Steinhausa. Konstrukcje tego typu
wiążą się z niepowodzeniem własności aproksymacji,
którą omówimy w następnym rozdziale. Istnienie
\emph{ciągu bazowego} w każdej nieskończenie wymiarowej
przestrzeni Banacha nie oznacza, że jest on bazą całej
przestrzeni.
\end{uwaga}
